\chapter{Giữ và Buông}
\markboth{Giữ và Buông}{Holding and Letting Go}

\section{Giữ một nguyên tắc giữa đám đông dễ dãi.}
\begin{truyen}{Giữ một nguyên tắc giữa đám đông dễ dãi.}{Some values must be held.}
\chuhoa{K}{hi cả nhóm rủ nhau chép bài, chỉ một bạn từ chối. Bạn ấy biết làm vậy sẽ bị chê là cứng nhắc, nhưng vẫn chọn giữ nguyên tắc học thật. Có những điều nếu buông ra cho dễ hòa nhập thì sau đó rất khó nhặt lại bằng lòng tự trọng.}
\textit{(When the whole group planned to cheat, only one student refused. That student knew the choice would invite mockery, yet still held to honest learning. Some things, once dropped for easier belonging, are very hard to recover in self-respect.)}
\end{truyen}
\begin{baihoc}
Không phải cái gì cũng nên mềm; có giá trị phải giữ đến cùng.
\textit{(Not everything should be softened; some values must be held to the end.)}
\end{baihoc}
\begin{ghinhoanh}
\textbf{Short English to remember:}

Some values must be held.
\end{ghinhoanh}
\begin{tuhocnhanh}
1. Nguyên tắc nào của em đang dễ bị đám đông kéo lệch? \textit{(Which of your principles is most vulnerable to crowd pressure?)}

2. Em có điều gì cần giữ rõ hơn trong thời gian này? \textit{(What do you need to hold more clearly these days?)}
\end{tuhocnhanh}
\ngancach

\section{Buông một kỳ vọng quá sức để sống thật hơn.}
\begin{truyen}{Buông một kỳ vọng quá sức để sống thật hơn.}{Letting go can make room for truth.}
\chuhoa{M}{ẹ luôn muốn con phải đứng đầu lớp, còn đứa con ngày một mệt và sợ học. Chỉ đến khi cả nhà chấp nhận mục tiêu vừa sức hơn, không khí mới dịu lại và việc học trở về đúng nghĩa của nó. Đôi lúc buông bớt kỳ vọng không làm con người kém đi, mà làm họ thở được.}
\textit{(A mother always wanted her child to rank first, while the child grew more exhausted and fearful of learning. Only when the family accepted a more realistic goal did the atmosphere soften and learning return to its true meaning. At times, letting go of excess expectation does not diminish a person; it allows breathing.)}
\end{truyen}
\begin{baihoc}
Buông đúng thứ có thể cứu cả một mối quan hệ.
\textit{(Letting go of the right thing can rescue an entire relationship.)}
\end{baihoc}
\begin{ghinhoanh}
\textbf{Short English to remember:}

Letting go can save love.
\end{ghinhoanh}
\begin{tuhocnhanh}
1. Kỳ vọng nào đang làm em hoặc người thân mệt quá mức? \textit{(Which expectation is exhausting you or your loved ones?)}

2. Nếu buông bớt một chút, điều gì sẽ lành lại? \textit{(If you let go a little, what might begin to heal?)}
\end{tuhocnhanh}
\ngancach

\section{Giữ ký ức đẹp nhưng buông nỗi cay cú.}
\begin{truyen}{Giữ ký ức đẹp nhưng buông nỗi cay cú.}{Keep the lesson, release the poison.}
\chuhoa{S}{au một tình bạn rạn vỡ, cô vẫn nhớ những năm tháng chân thành đã có. Nhưng cô không muốn giữ mãi cơn giận vì biết nó sẽ làm mình đắng đi. Trưởng thành không phải xóa hết chuyện cũ, mà là biết giữ điều đáng quý và buông điều đang đầu độc mình.}
\textit{(After a friendship broke, she still remembered the sincere years they had shared. Yet she did not want to preserve the bitterness because she knew it would make her hard inside. Maturity is not erasing the past, but keeping what is worthy and releasing what is poisoning you.)}
\end{truyen}
\begin{baihoc}
Giữ và buông không đối chọi nhau; chúng cùng phục vụ cho sự lành mạnh của tâm hồn.
\textit{(Holding and letting go are not enemies; both serve the health of the soul.)}
\end{baihoc}
\begin{ghinhoanh}
\textbf{Short English to remember:}

Keep the lesson, release the poison.
\end{ghinhoanh}
\begin{tuhocnhanh}
1. Điều gì trong quá khứ em nên giữ lại như một bài học? \textit{(What from your past should you keep as a lesson?)}

2. Điều gì em cần buông vì nó chỉ làm lòng mình nặng thêm? \textit{(What do you need to release because it only burdens your heart?)}
\end{tuhocnhanh}
\ngancach

\section{Buông sĩ diện để học điều mình chưa biết.}
\begin{truyen}{Buông sĩ diện để học điều mình chưa biết.}{Pride dropped becomes space for learning.}
\chuhoa{A}{nh đã đi làm nhiều năm nên ngại hỏi những điều căn bản trước đồng nghiệp trẻ hơn. Càng ngại, anh càng mắc lỗi. Đến khi dám buông sĩ diện và hỏi thật, công việc lại nhẹ đi rất nhiều.}
\textit{(He had worked for years and felt embarrassed asking basic questions before younger colleagues. The more he hid, the more mistakes appeared. Once he let go of pride and honestly asked, the work became much lighter.)}
\end{truyen}
\begin{baihoc}
Buông cái tôi không làm ta nhỏ đi; nó làm ta học được.
\textit{(Letting go of ego does not make us smaller; it makes learning possible.)}
\end{baihoc}
\begin{ghinhoanh}
\textbf{Short English to remember:}

Dropped pride opens learning.
\end{ghinhoanh}
\begin{tuhocnhanh}
1. Em đang giấu sự chưa biết ở chỗ nào chỉ vì ngại? \textit{(Where are you hiding your lack of knowledge out of embarrassment?)}

2. Nếu dám hỏi một câu, em sẽ nhẹ đi điều gì? \textit{(If you dared to ask one question, what burden would become lighter?)}
\end{tuhocnhanh}
\ngancach

\section{Giữ mình trong sạch trước một lợi ích quá dễ.}
\begin{truyen}{Giữ mình trong sạch trước một lợi ích quá dễ.}{What is easy is not always worth keeping.}
\chuhoa{C}{ó lúc cơ hội đến rất nhanh: chỉ cần nhắm mắt một chút, im lặng một chút, thỏa hiệp một chút là được lợi. Nhưng anh biết nếu nhận cái dễ đó, mình sẽ phải trả bằng một điều khó hơn nhiều: sự xem nhẹ chính mình. Có cái nên buông ngay từ đầu để còn giữ được điều quan trọng hơn.}
\textit{(Sometimes opportunity arrives very quickly: close your eyes a little, stay silent a little, compromise a little, and gain follows. Yet he knew that taking such ease would require paying a far heavier price: losing respect for himself. Some things must be let go early in order to keep what matters more.)}
\end{truyen}
\begin{baihoc}
Giữ mình nhiều khi bắt đầu bằng việc buông một món lợi không sạch.
\textit{(Keeping yourself often begins by letting go of an unclean advantage.)}
\end{baihoc}
\begin{ghinhoanh}
\textbf{Short English to remember:}

Let go of the easy wrong.
\end{ghinhoanh}
\begin{tuhocnhanh}
1. Món lợi nào dễ quá đang làm em nên cẩn thận hơn? \textit{(Which too-easy gain should make you more cautious?)}

2. Em muốn giữ điều gì lâu dài hơn một lợi ích trước mắt? \textit{(What do you want to preserve that matters more than immediate benefit?)}
\end{tuhocnhanh}
\ngancach
